%% --------------------------------------------------------
\section{Порівняння CI та CC методів}
%% --------------------------------------------------------

\inputcode{CIvsCD.py}

%% --------------------------------------------------------
\subsection{Size-extensivity проблема CI}
%% --------------------------------------------------------

Важлива відмінність між CI та CC:

\begin{equation}
    E_{CI}(2A) \neq 2 \times E_{CI}(A)
\end{equation}

\begin{equation}
    E_{CC}(2A) = 2 \times E_{CC}(A)
\end{equation}

\inputcode{SizeExtensivityExample.py}


%% --------------------------------------------------------
\subsection{Коли використовувати CI vs CC}
%% --------------------------------------------------------
\nopagebreak
\begin{center}
    \captionof{table}{Порівняння CI та CC методів}
    \label{tab:ci_vs_cc}
    \small
    \begin{tblr}{lll}
        \hline
        \textbf{Критерій}       & \textbf{CI}      & \textbf{CC}      \\
        \hline
        Варіаційність           & Так              & Ні               \\
        Size-extensivity        & Ні (окрім FCI)   & Так              \\
        Точність                & CISD < CCSD      & CCSD(T) ≈ FCI    \\
        Збуджені стани          & EOM-CI, CIS      & EOM-CC           \\
        Швидкість               & CISD швидше CCSD & Повільніше       \\
        Стабільність            & Варіаційна межа  & Може розходитись \\
        Багатоконфігураційність & Природно (FCI)   & Потребує MRCC    \\
        \hline
    \end{tblr}
\end{center}

\textbf{Рекомендації:}
\begin{itemize}
    \item \textbf{Використовуйте FCI} для малих систем як еталон
    \item \textbf{Використовуйте CCSD(T)} для точних розрахунків одноконфігураційних систем
    \item \textbf{Використовуйте CISD} для швидких оцінок кореляції
    \item \textbf{Для багатоконфігураційних систем} використовуйте CASSCF/CASPT2
\end{itemize}